\usepackage{iftex} % Добавляем пакет для определения используемого движка TeX

%%% Проверка используемого TeX-движка %%%
\newif\ifxetexorluatex   % определяем новый условный оператор
\ifXeTeX
    \xetexorluatextrue
\else
    \ifLuaTeX
        \xetexorluatextrue
    \else
        \xetexorluatexfalse
    \fi
\fi

\usepackage{etoolbox}   % Для продвинутой проверки разных условий
\providebool{presentation}

\usepackage{comment}    % Позволяет убирать блоки текста

%%% Поля и разметка страницы %%%
\usepackage{pdflscape}  % Для включения альбомных страниц
\usepackage{geometry}   % Для последующего задания полей

%%% Форматирование текста %%%
\usepackage{paracol}

%%% Математические пакеты %%%
\usepackage{amsthm,amsmath,amscd}   % Математические дополнения от AMS
\usepackage{amsfonts,amssymb}       % Математические дополнения от AMS
\usepackage{mathtools}              % Добавляет окружение multlined
\usepackage{xfrac}                  % Красивые дроби
\usepackage[
    locale = DE,
    list-separator       = {;\,},
    list-final-separator = {;\,},
    list-pair-separator  = {;\,},
    list-units           = single,
    range-units          = single,
    range-phrase={\text{\ensuremath{-}}},
    fraction-function    = \sfrac,
    separate-uncertainty,
    ]{siunitx}                 % Размерности SI
\sisetup{inter-unit-product = \ensuremath{{}\cdot{}}}

% Кириллица в нумерации subequations
\usepackage{alphalph} % Для корректной работы \asbuk
\usepackage{chngcntr} % Для сброса счётчиков

\patchcmd{\subequations}{\def\theequation{\theparentequation\alph{equation}}}
{\def\theequation{\theparentequation\asbuk{equation}}}
{\typeout{subequations patched}}{\typeout{subequations not patched}}

%%%% Установки для размера шрифта 14 pt %%%%
%% Формирование переменных и констант для сравнения
\newlength{\curtextsize}
\newlength{\bigtextsize}
\setlength{\bigtextsize}{13.9pt}

\makeatletter
\setlength{\curtextsize}{\f@size pt}
\makeatother

%%% Кодировки и шрифты %%%
\ifxetexorluatex
    \ifpresentation
        \providecommand*\autodot{} % quick fix for polyglossia
    \fi
    \PassOptionsToPackage{no-math}{fontspec}
    \usepackage{polyglossia}
    \setdefaultlanguage{russian}
    \setotherlanguage{english}
\else
    \usepackage{cmap}   % Улучшенный поиск русских слов в PDF
    \usepackage{textcomp}
    \usepackage[english, greek, russian]{babel}
    \usepackage[T1,T2A]{fontenc}
    \usepackage[utf8]{inputenc}
    \makeatletter\AtBeginDocument{\let\@elt\relax}\makeatother
\fi

%%% Оформление абзацев %%%
\ifpresentation
\else
    \usepackage{indentfirst} % Красная строка в начале главы
\fi

%%% Оформление цитат %%%
\usepackage{csquotes}

%%% Оформление эпиграфов %%%
\usepackage{epigraph}

%%% Цвета %%%
\ifpresentation
\else
    \usepackage[dvipsnames, table, hyperref]{xcolor}
\fi

%%% Таблицы %%%
\usepackage{longtable,ltcaption}
\usepackage{multirow,makecell}
\usepackage{tabularx} % Заменяем tabu на tabularx
\usepackage{array} % Для настраиваемых типов колонок
\usepackage{ragged2e} % Для управления выравниванием текста

\usepackage{threeparttable}

%%% Общее форматирование %%%
\usepackage{soul}   % Заменяем soulutf8 на soul
\usepackage{icomma}

%%% Векторная графика %%%
\usepackage{tikz}
\usetikzlibrary{chains, shapes.geometric, shapes.symbols, arrows}
\usetikzlibrary{patterns}

\usepackage{pgfplots} % Оформление графиков
\usepackage{pgfplotstable}
\usepgfplotslibrary{dateplot}
\usepgfplotslibrary{fillbetween}
\usepgfplotslibrary{statistics}
\pgfplotsset{compat=newest} % Allows to place the legend below plot
\usepgfplotslibrary{units} % Allows to enter the units nicely
\sisetup{
  round-mode = places,
  round-precision = 2,
}

%%% Гиперссылки %%%
\ifxetexorluatex
    \let\CYRDZE\relax
\fi
\usepackage{hyperref}

%%% Изображения %%%
\usepackage{graphicx}
\usepackage{caption}
\usepackage{subcaption}
\usepackage{pdfpages}

%%% Счётчики %%%
\usepackage{aliascnt}
\usepackage[figure,table]{totalcount}
\usepackage{totcount}
\usepackage{totpages}

%%% Продвинутое управление ссылками %%%
\ifpresentation
\else
    \usepackage{cleveref}
    \usepackage{kvsetkeys}
    \makeatletter
    \let\org@@cref\@cref
    \renewcommand*{\@cref}[2]{%
        \edef\process@me{%
            \noexpand\org@@cref{#1}{\zap@space#2 \@empty}%
        }\process@me
    }
    \makeatother
\fi

\usepackage{placeins}

\ifnumequal{\value{draft}}{1}{% Черновик
    \usepackage[firstpage]{draftwatermark}
    \SetWatermarkText{DRAFT}
    \SetWatermarkFontSize{14pt}
    \SetWatermarkScale{15}
    \SetWatermarkAngle{45}
}{}

%%% Списки %%%
\usepackage{enumitem}

%%% Оформление списка обозначений %%%
\usepackage[intoc]{nomencl}
\makenomenclature
\setlength{\nomitemsep}{-.8\parsep}

\usepackage{fr-longtable}    % Для \endlasthead

% Листинги исходного кода
\usepackage{fancyvrb}
\usepackage{listings}
\lccode`\~=0\relax % Исправление для listings в XeLaTeX и LuaLaTeX

% Русская традиция начертания греческих букв
\usepackage{upgreek}

%%% Отметка о версии черновика на каждой странице %%%
\ifnumequal{\value{draft}}{1}{% Черновик
    \IfFileExists{.git/gitHeadInfo.gin}{
        \usepackage[mark,pcount]{gitinfo2}
        \renewcommand{\gitMark}{rev.\gitAbbrevHash\quad\gitCommitterEmail\quad\gitAuthorIsoDate}
        \renewcommand{\gitMarkFormat}{\rmfamily\color{Gray}\small\bfseries}
    }{}
}{}
